Note: for questions requiring giving a position, correct to $1^\circ$ gets
full marks (usually 1 point), correct to $2^\circ$ gets half marks (usually
0.5 points).

\begin{description}
\item[(A)] The nova, the Moon and the comet are marked on the sky map (see
  the marked answer map). \hfill(2+1+3 points)
  % FIGURE NEEDED: marked answer sky map, 2011_PLAN_sol.pdf page 2 (polar
  % projection with N = nova, Moon symbol, comet shape, Z + local meridian
  % line, M = Martian pole, P/PE = ecliptic north pole, L = red-arrow star)
\item[(B)] Objects above the astronomical horizon (to be circled):
  $o$ Cet -- Mira; $\delta$ CMa -- Wezen; $\alpha$ Cyg -- Deneb;
  $\beta$ Per -- Algol; $\delta$ Cep -- Alrediph; M44 -- Praesepe (Beehive
  Cluster). Below the horizon (not circled): M20 -- Trifid Nebula;
  M57 -- Ring Nebula; $\alpha$ Boo -- Arcturus.
  \hfill(+0.5 for each correct, $-1$ for each incorrect)
\item[(C)] The northern part of the local meridian and the ecliptic north
  pole are marked on the answer map. \hfill(3+2 points)
\item[(D)] Geographical latitude of the observer:
  $\varphi = 50^\circ \pm 1^\circ$; \hfill(1 point)\\
  Local Sidereal Time:
  $\theta = 5^{\mathrm{h}}25^{\mathrm{m}} \pm 8^{\mathrm{m}}$;
  \hfill(2 points)\\
  time of year: Jan or Feb. \hfill(2 points)
\item[(E)] $A_1 = 45^\circ$, $h_1 = 58^\circ$: M45; \hfill(1 point)\\
  $A_2 = 278^\circ$, $h_2 = 20^\circ$: $\alpha$ Leo. \hfill(1 point)
\item[(F)] Sirius ($\alpha$ CMa):
  $A_3 = 340^\circ \pm 2^\circ$, $h_3 = 20^\circ \pm 1^\circ$;
  \hfill(0.5+0.5 points)\\
  the Andromeda Galaxy (M31):
  $A_4 = 108^\circ \pm 2^\circ$, $h_4 = 41^\circ \pm 2^\circ$.
  \hfill(0.5+0.5 points)
\item[(G)] $\alpha = 3^{\mathrm{h}}24^{\mathrm{m}} \pm 15^{\mathrm{m}}$;
  $\delta = 49^\circ \pm 3^\circ$. \hfill(1+1 points)
\end{description}
